% xelatex
\documentclass[11pt,a4paper]{article}
\usepackage{xeCJK}
\usepackage{geometry}
\usepackage{amsmath,amssymb,amsfonts,amsthm}
\usepackage[dvipsnames,svgnames]{xcolor}
\usepackage{enumerate,enumitem}
\usepackage{thmtools}
\usepackage{fancyhdr}
\usepackage{ragged2e}
\usepackage{titling}
\usepackage{hyperref}
\usepackage{graphicx}
\usepackage{float}
\usepackage{subcaption}


% ---------- Preamble ----------
\pagestyle{fancyplain} % Make all pages in the document conform to the custom headers and footers
\fancyhead{} % No page header
\fancyfoot[L]{}
\fancyfoot[C]{}
\fancyfoot[R]{\thepage}
\renewcommand{\headrulewidth}{0pt} % Remove header underlines
\renewcommand{\footrulewidth}{0pt} % Remove footer underlines
\geometry{
    textwidth=160mm,
    top=25mm,
    bottom=25mm,
}
\setlist[enumerate]{
    label=\textnormal{(\arabic*)},
    topsep=.5em,
    itemsep=0em,
}
\setlist[itemize]{
    topsep=.5em,
    itemsep=0em,
}


% ---------- Theorem environments ----------
\theoremstyle{plain} % default
\newtheorem{thm}{Theorem}
\newtheorem*{theorem}{Theorem}
\newtheorem{lemma}[thm]{Lemma}

\theoremstyle{definition}
\newtheorem*{example}{Example}

\theoremstyle{remark}
\newtheorem*{rmk}{Remark}


% ---------- Symbol Macros ----------
\newcommand{\bra}[1]{\mathopen{}\left(#1\right)}
\newcommand{\sbra}[1]{\mathopen{}\left[#1\right]}
\newcommand{\cbra}[1]{\mathopen{}\left\{#1\right\}}
\newcommand{\abra}[1]{\mathopen{}\left\langle#1\right\rangle}
\newcommand{\norm}[1]{\mathopen{}\left\lVert#1\right\rVert}
\newcommand{\inp}[2]{\mathopen{}\left\langle#1,#2\right\rangle}
\newcommand{\abs}[1]{\mathopen{}\left|#1\right|}
\newcommand{\rest}[2]{\mathopen{}\left.#1\right|_{#2}}

\newcommand{\R}{\mathbf{R}}
\newcommand{\N}{\mathbf{N}}
\newcommand{\Z}{\mathbf{Z}}
\newcommand{\C}{\mathbf{C}}
\newcommand{\Q}{\mathbf{Q}}

\renewcommand{\d}{\mathrm{d}}
\newcommand{\dx}{\mathrm{d}x}
\newcommand{\eps}{\varepsilon}
% \renewcommand{\phi}{\varphi}
\renewcommand{\L}{\mathcal{L}}
\renewcommand{\S}{\mathcal{S}}

\newcommand{\nti}{n\to\infty}


% ---------- Title Section ----------
\title{A Natural Proof of the Dominated Convergence Theorem}
\author{Zhijie Chen}

\begin{document}
\maketitle

The Dominated Convergence Theorem is a grand result in the area of Lebesgue integration. It deals with the problem of interchanging limits and integrals, i.e., whether
\[\lim_{\nti}\int f_n=\int f\]
for $f_n\to f$. The theorem demonstrates that Lebesgue integration indeed fixes a serious deficiency of Riemann integration, at least to a reasonably satisfactory extent.

\begin{theorem}[Dominated Convergence Theorem, DCT]
    Suppose $(X,\S,\mu)$ is a measure space, $f:X\to[-\infty,\infty]$ is measurable, and $\{f_n\}_{n\in\N}$ is a sequence of measurable functions such that $f_n\to f$ a.e. If there exists a nonnegative measurable function $g\in\L^1(X,\S,\mu)$ s.t.
    \[\abs{f_n(x)}\le g(x)\]
    for every $n\in\N$ and almost every $x\in X$, then
    \[\lim_{\nti}\int f_n\,\d\mu=\int f\,\d\mu.\]
\end{theorem}

The traditional proof of DCT uses Fatou's Lemma and is sketched in \ref{subsec:proof_dct_fatou}. However, the proof is clean but not natural, at least to me. This article presents an alternative proof using Egorov's Theorem, which is slightly more complicated but quite intuitive.

\section{Background}
For the rest of the article, suppose $(X,\S,\mu)$ is a measure space.

\subsection{A deficiency of Riemann integration}
One problem with Riemann integration is that it does not work well with limits. Integrals of (pointwise) limits do not necessarily equal limits of integrals; limits of Riemann integrable functions are not necessarily Riemann integrable\footnote{In contrast, the pointwise limit of measurable functions is guaranteed to be measurable.}. Below are two interesting examples.

\begin{example}
    Define $f_n:[0,1]\to\R$ for each $n\in\N$ by
    \[f_n(x)=\lim_{m\to\infty}\abs{\cos(2n!\pi x)}^m.\]
    Then $\int f_n=0$ for each $n$ but
    \[f(x)=\lim_{\nti}f_n(x)=\chi_\Q\]
    is the Dirichlet function and is not Riemann integrable.
\end{example}

\begin{example}[Example 1.16 of \emph{Measure, Integration, and Real Analysis} by Sheldon Axler]
    Suppose $\{r_n\}_{n\in\N}$ is an enumeration of $(0,1)\cap\Q$. Define $f_n:[0,1]\to\R$ for each $n\in\N$ by
    \[f_n(x)=\begin{cases}
        \frac1{\sqrt{x-r_n}}\quad&\text{if }x>r_n,\\
        0&\text{otherwise}.
    \end{cases}\]
    Define $f:[0,1]\to[0,\infty]$ by
    \[f(x)=\sum_{n=0}^{\infty}\frac{f_n(x)}{2^n}.\]
    Because $\Q$ is dense in $\R$, $f$ is unbounded on every nontrivial interval. Thus $f$ is not Riemann integrable on every nontrivial subinterval of $[0,1]$. However, we should expect the area under the graph of $f$ to be less than 4 rather than undefined, since the area under the graph of each $f_n$ is less than 2.
\end{example}

Although we have the following positive result, the theorem has two problems: it requires the hypothesis that the limit function is Riemann integrable, and mathematicians have not yet found a clean, elementary, Riemann-integration-based proof of it. In fact, this suggests that Riemann integration is not the ideal theory of integration. See Theorem~1.18 of \emph{Measure, Integration, and Real Analysis} by Sheldon Axler and section~14.2.3 of 《数学分析习题课讲义（第 2 版）（下册）》 by 谢惠民 et al.

\begin{theorem}[Arzelà's Bounded Convergence Theorem]
    Suppose $\{f_n\}_{n\in\N}$ is a sequence of Riemann integrable functions on $[a,b]$, $\{f_n\}$ is uniformly bounded, and $f_n\to f$. If $f$ is Riemann integrable, then
    \[\lim_{\nti}\int_a^bf_n=\int_a^bf.\]
\end{theorem}

\subsection{Traditional proof of DCT by Fatou's Lemma}
\label{subsec:proof_dct_fatou}
We state Fatou's Lemma and sketch the proof of DCT using Fatou's Lemma below.

\begin{theorem}[Fatou's Lemma]
    Suppose $\{f_n\}_{n\in\N}$ is a sequence of nonnegative measurable functions. Then
    \[\int\bra{\liminf_{\nti}f_n}\d\mu\le\liminf_{\nti}\int f_n\,\d\mu.\]
\end{theorem}

\begin{proof}[Proof sketch of DCT using Fatou's Lemma]
    Notice that $g-f_n\ge0$ a.e. and $g+f_n\ge0$ a.e. By Fatou's Lemma,
    \[\int\bra{\liminf_{\nti}(g-f_n)}\d\mu\le\liminf_{\nti}\int (g-f_n)\,\d\mu \implies \int f\,\d\mu\ge\limsup_{\nti}\int f_n\,\d\mu,\]
    \[\int\bra{\liminf_{\nti}(g+f_n)}\d\mu\le\liminf_{\nti}\int (g+f_n)\,\d\mu \implies \int f\,\d\mu\le\liminf_{\nti}\int f_n\,\d\mu.\]
    The proof is completed.
\end{proof}

\section{Some Preparation}
We first state Egorov's Theorem.

\begin{theorem}[Egorov's Theorem]
    Suppose $\mu(X)<\infty$, $\{f_n\}_{n\in\N}$ is a sequence of measurable functions, and $f_n\to f$. Then $\forall\eps>0$, $\exists E\in\S$ s.t. $\mu(E^c)<\eps$ and $f_n$ converges uniformly to $f$ on $E$.
\end{theorem}

To use Egorov's Theorem to prove DCT, we still need to address two problems.
\begin{enumerate}
    \item $\mu(X)$ might not be finite.
    \item We need to control the integral over the set on which $f_n$ does not converge uniformly.
\end{enumerate}
Lemma~\ref{lem:mass_concentrate} will solve the first problem, and Lemma~\ref{lem:small_area} will solve the second.

\subsection{Integrals over small sets are small}
\begin{lemma}\label{lem:small_area}
    Suppose $f\in\L^1(X,\S,\mu)$ is nonnegative. Then $\forall\eps>0$, $\exists\delta>0$ s.t. for every $B\in\S$ with $\mu(B)<\delta$,
    \[\int_Bf\,\d\mu<\eps.\]
\end{lemma}

\begin{proof}
    Suppose $\eps>0$. The desired result is trivial for simple functions. For the general case, take a simple measurable function $h$ s.t. $0\le h\le f$ and $\int(f-h)\,\d\mu<\eps/2$. Then $\exists\delta>0$ s.t. for every $B\in\S$ with $\mu(B)<\delta$,
    \[\int_Bh\,\d\mu<\frac\eps2.\]
    Then
    \[\int_Bf\,\d\mu\le\int(f-h)\,\d\mu+\int_Bh\,\d\mu<\eps\]
    for every $B\in\S$ with $\mu(B)<\delta$, as desired.
\end{proof}

\subsection{Mass of integrable functions concentrates on sets of finite measure}
\begin{lemma}\label{lem:mass_concentrate}
    Suppose $f\in\L^1(X,\S,\mu)$ is nonnegative. Then $\forall\eps>0$, $\exists E\in\S$ s.t. $\mu(E)<\infty$ and
    \[\int_{E^c}f\,\d\mu<\eps.\]
\end{lemma}
\begin{proof}
    Suppose $\eps>0$. The desired result is trivial for simple functions. For the general case, take a simple measurable function $h$ s.t. $0\le h\le f$ and $\int(f-h)\,\d\mu<\eps/2$. Then $\exists E\in\S$ s.t. $\mu(E)<\infty$ and
    \[\int_{E^c}h\,\d\mu<\frac\eps2.\]
    Then
    \[\int_{E^c}f\,\d\mu\le\int(f-h)\,\d\mu+\int_{E^c}h\,\d\mu<\eps,\]
    as desired.
\end{proof}

\section{A Natural Proof of DCT by Egorov's Theorem}
Fix $\eps>0$. By Lemma~\ref{lem:mass_concentrate}, let $X_0\in\S$ be s.t. $\mu(X_0)<\infty$ and
\[\int_{X_0^c}g\,\d\mu<\frac\eps5.\]
By Lemma~\ref{lem:small_area}, let $\delta>0$ be s.t for every $B\in\S$ with $\mu(B)<\delta$,
\[\int_Bg\,\d\mu<\frac\eps5.\]
By Egorov's Theorem, $\exists E\subseteq X_0$ s.t. $E\in\S$, $\mu(X_0\setminus E)<\delta$, and $f_n$ converges uniformly to $f$ on $E$. Then
\begin{align*}
    \abs{\int(f_n-f)\,\d\mu}&\le\int\abs{f_n-f}\,\d\mu \\
    &\le\int_E\abs{f_n-f}\,\d\mu + \int_{X_0\setminus E}(\abs{f_n}+\abs f)\,\d\mu + \int_{X_0^c}(\abs{f_n}+\abs f)\,\d\mu \\
    &\le\int_E\abs{f_n-f}\,\d\mu + 2\int_{X_0\setminus E}g\,\d\mu + 2\int_{X_0^c}g\,\d\mu \\
    (n\ge N)\quad&\le\eps,
\end{align*}
where $N\in\N$ is s.t.
\[\abs{f_n(x)-f(x)}<\frac\eps{5\mu(E)}\]
for every $n\ge N$ and every $x\in E$. The proof is completed.\qed

\end{document}
